% -----------------------------------------------------------
\section{Spirit}
% -----------------------------------------------------------
	\label{SPIRIT}
	
% % ----------------------
% \subsection{See also}
% % ----------------------

% ----------------------
\subsection{1828 American Dictionary of the English Language} % *1828
% ----------------------
	\label{SPIRIT-1828}

	\begin{compactitem}
		\item[11] Eager desire; disposition of mind excited and directed to a particular object. God has made a \textit{spirit} of building succeed a \textit{spirit} of pulling down.
		\item[15] Life or strength of resemblance; essential qualities; as, to set off the face in its true \textit{spirit} The copy has not the \textit{spirit} of the original.
		\item[17] That which hath power or energy; the quality of any substance which manifest life, activity, or the power of strongly affecting other bodies; as the \textit{spirit} of wine or of any liquor.
	\end{compactitem}

% % ----------------------
% \subsection{Oxford English Dictionary} % *OED
% % ----------------------
	% \label{SPIRIT-OED}

	% \begin{compactitem}
		% \item[]
	% \end{compactitem}

% ----------------------
\subsection{Scriptural Usage} % *SCRIPTURES
% ----------------------
	\label{SPIRIT-SCRIPTURES}

	% % -----
	% \subsubsection{Old Testament} % *OT
	% % -----
		% \label{SPIRIT-OT}
	
		% % --
		% \paragraph{KJV}
		% % --
			% \label{SPIRIT-OT-KJV}
	
			% \begin{tabular}{ll}
				% Total occurrences in the OT & 
			% \end{tabular}
			
			% \begin{compactdesc}
				% \item[]
			% \end{compactdesc}
			
		% % --
		% \paragraph{Hebrew Bible}
		% % --
			% \label{SPIRIT-HEBREW}
		
			% %
			% \subparagraph{\texorpdfstring{\he{sample word}}{}}
			% %
				% \label{PAGA} % whatever is in step bible without periods
			
				% \begin{compactdesc}
					% \item[HALOT , PDF ] \textbf{}
						% \begin{compactitem}
							% \item[1]
						% \end{compactitem}
					% \item[BDB , PDF ] \textbf{}
						% \begin{compactitem}
							% \item[1]
						% \end{compactitem}
					% \item[DCH PDF ] \textbf{}
						% \begin{compactitem}
							% \item[1]
						% \end{compactitem}
				% \end{compactdesc}
				
				% \begin{compactdesc}
					% \item[Some OT uses] \textbf{}
						% \begin{compactdesc}
							% \item[]
						% \end{compactdesc}
				% \end{compactdesc}
			
		% % --
		% \paragraph{Septuagint}
		% % --
			% \label{SPIRIT-LXX}
		
			% %
			% \subparagraph{\texorpdfstring{\gr{sample word}}{}}
			% %
		
	% -----
	\subsubsection{New Testament} % *NT
	% -----
		\label{SPIRIT-NT}
	
		% % --
		% \paragraph{KJV}
		% % --
			% \label{SPIRIT-NT-KJV}
			
	
			% \begin{tabular}{ll}
				% Total occurrences in the NT & 
			% \end{tabular}
			
			% \begin{compactdesc}
				% \item[]
			% \end{compactdesc}
			
		% --
		\paragraph{Greek New Testament} % *GREEK
		% --
			\label{SPIRIT-GREEK}
		
			%
			\subparagraph{\texorpdfstring{\gr{πνεῦμα}}{}}
			%
				\label{PNEUMA}
			
				\begin{compactdesc}
					\item[BDAG 738b] \textbf{}
						\begin{compactitem}
							\item[1] air in movement, \textit{blowing, breathing}
							\item[1a] wind
							\item[1b] the breathing out of air, blowing, breath
							\item[2] that which animates or gives life to the body, \textit{breath, (life-)spirit}
							\item[3] a part of human personality, \textit{spirit}
							\item[3a] when used with \gr{σάρξ}, the flesh, it denotes the immaterial part...flesh and spirit = the whole personality, in its outer and inner aspects \textit{(There is more interesting material here.)}
							\item[3b] as the source and seat of insight, feeling, and will, generally as the representative part of human inner life
							\item[3c] spiritual state, state of mind, disposition
							\item[4] an independent noncorporeal being, in contrast to a being that can be perceived by the physical senses, \textit{spirit}
							\item[4a] God personally
							\item[4b] good, or at least not expressly evil \textit{spirits} or \textit{spirit-beings}
							\item[4c] evil \textit{spirits}
							\item[5] God's being as controlling influence, with focus on association with humans, \textit{Spirit, spirit} as that which differentiates God from everything that is not God, as the divine power that produces all divine existence, as the divine element in which all divine life is carried on, as the bearer of every application of the divine will. All those who belong to God possess or receive this spirit and hence have a share in God's life. This spirit also serves to distinguish Christians from all unbelievers
							\item[5a] the Spirit of God, of the Lord (=God)
							\item[5b] the Spirit of Christ, of the Lord (=Christ)
							\item[5c] Because of its heavenly origin and nature this Spirit is called \textit{(the) Holy Spirit}
							\item[5c$\alpha$] with the article 
							\item[5c$\beta$] without the article
							\item[5d] absolute
							\item[5d$\alpha$] with the article \gr{τό πνεῦμα}
							\item[5d$\beta$] without the article \gr{πνεῦμα}
							\item[5e] The Spirit is more closely defined by a genitive of thing
							\item[5f] Of Christ ``it is written'' in Scripture: \gr{Ἐγένετο...ὁ ἔσχατος Ἀδὰμ εἰς πνεῦμα ζῳοποιοῦν} (1 Corinthians 15:45).
							\item[5g] The (divine) Pneuma stands in contrast to everything that characterizes this age or the finite world generally
							\item[5g$\alpha$] in contrast to \gr{σάρξ}, which is more closely connected with sin than any other earthly material
							\item[5g$\beta$] in contrast to \gr{σῶμα}
							\item[5g$\gamma$] in contrast to \gr{γράμμα}, which is the characteristic quality of God's older declaration of the divine will in the law
							\item[5g$\delta$] in contrast to the wisdom of humans
							\item[6] the Spirit of God as exhibited in the character or activity of God's people or selected agents, \textit{Spirit, spirit}
							\item[6a] \gr{πνεῦμα} is accompanied by another noun, which characterizies the working of the Spirit more directly
							\item[6b] unless frustrated by humans in their natural condition, the Spirit of God produces a spiritual type of conduct
							\item[6c] The Spirit inspires certain people of God, above all, in their capacity as proclaimers of a divine revelation
							\item[6d] The Spirit of God, beineg one, shows the variety and richness of its life in the different kinds of spiritual gifts which are granted to certain Christians
							\item[6e] One special type of spiritual gift is represented by ecstatic speaking. Of those who ``speak in tongues'' that no earthly person can understand
							\item[6f] The Spirit leads and directs Christian missionaries in their journeys
							\item[7] an activating spirit that is not from God, \textit{spirit}
							\item[8] an independent transcendent personality, \textit{the Spirit}, which appears in formulas that became more and more fixed and distinct \textit{(Much more available here, some description but mostly citing outside research.)}
						\end{compactitem}
					\item[LSJ , PDF ] \textbf{}
						\begin{compactitem}
							\item 
						\end{compactitem}
				\end{compactdesc}
				
				\begin{compactdesc}
					\item[Some NT uses] \textbf{}
						\begin{compactdesc}
							\item[]
						\end{compactdesc}
				\end{compactdesc}
		
	% % -----
	% \subsubsection{Book of Mormon} % *BOM
	% % -----
		% \label{SPIRIT-BOM}
	
		% \begin{tabular}{ll}
			% Total occurrences in the BOM & 
		% \end{tabular}
		
		% \begin{compactdesc}
			% \item[]
		% \end{compactdesc}
		
	% % -----
	% \subsubsection{Doctrine and Covenants} % *DC
	% % -----
		% \label{SPIRIT-DC}
	
		% \begin{tabular}{ll}
			% Total occurrences in the DC & 
		% \end{tabular}
		
		% \begin{compactdesc}
			% \item[]
		% \end{compactdesc}
		
	% % -----
	% \subsubsection{Pearl of Great Price} % *PGP
	% % -----
		% \label{SPIRIT-PGP}
	
		% \begin{tabular}{ll}
			% Total occurrences in the PGP & 
		% \end{tabular}
		
		% \begin{compactdesc}
			% \item[]
		% \end{compactdesc}
		
% ----------------------
\subsection{Key Phrases} % *KEYPHRASES *PHRASES
% ----------------------
	\label{SPIRIT-PHRASES}

	\begin{compactdesc}
	
		\item[contrite spirit] \textbf{20} | \hyperref[PSALMS-34-18]{Psalms 34:18}; \hyperref[PSALMS-51-17]{Psalms 51:17}; \hyperref[ISAIAH-57-15]{Isaiah 57:15}; \hyperref[ISAIAH-66-2]{Isaiah 66:2}; \hyperref[2NEPHI-2-7]{2 Nephi 2:7}; \hyperref[2NEPHI-4-32]{2 Nephi 4:32}; \hyperref[HELAMAN-8-15]{Helaman 8:15}; \hyperref[3NEPHI-9-20]{3 Nephi 9:20} (2); \hyperref[3NEPHI-12-19]{3 Nephi 12:19}; \hyperref[MORMON-2-14]{Mormon 2:14}; \hyperref[ETHER-4-15]{Ether 4:15}; \hyperref[MORONI-6-2]{Moroni 6:2}; \hyperref[DC20-37]{DC 20:37}; \hyperref[DC52-15]{DC 52:15}; \hyperref[DC52-16]{DC 52:16}; \hyperref[DC56-17]{DC 56:17}; \hyperref[DC56-18]{DC 56:18}; \hyperref[DC59-8]{DC 59:8}; \hyperref[DC97-8]{DC 97:8}
			\begin{compactdesc}
				\item[Bible anchor verses] \phantom{.}
					\begin{compactdesc}
						\item[{\hyperref[PSALMS-34-18]{Psalms 34:18}}]
						\item[{\hyperref[PSALMS-51-17]{Psalms 51:17}}]
						\item[{\hyperref[ISAIAH-57-15]{Isaiah 57:15}}]
						\item[{\hyperref[ISAIAH-66-2]{Isaiah 66:2}}]
					\end{compactdesc}
				% \item[Commentary] ---
			\end{compactdesc}
	
		\item[power of * spirit] See \hyperref[POWER-PHRASES]{here}.
		
		\item[spirit of Christ] \textbf{7} | \hyperref[ROMANS-8-9]{Romans 8:9}; \hyperref[PHILIPPIANS-1-19]{Philippians 1:19}; \hyperref[1PETER-1-11]{1 Peter 1:11}; \hyperref[MORONI-7-16]{Moroni 7:16}; \hyperref[MORONI-10-17]{10:17}; \hyperref[DC20-37]{DC 20:37}; \hyperref[DC84-45]{84:45--46}
			% \begin{compactdesc}
				% \item[Bible anchor verses] \phantom{.}
					% \begin{compactdesc}
						% \item[Romans 5:5] \gr{}
					% \end{compactdesc}
				% \item[Commentary] ---
			% \end{compactdesc}
			
		\item[spirit of prophecy] \textbf{21} | \hyperref[REVELATIONNT-19-10]{Revelation 19:10}
			\begin{compactdesc}
				\item[Bible anchor verses] \phantom{.}
					\begin{compactdesc}
						\item[Revelation 19:10] \gr{...ἡ γὰρ μαρτυρία Ἰησοῦ ἐστιν τὸ πνεῦμα τῆς προφητείας}
					\end{compactdesc}
				% \item[Commentary] ---
			\end{compactdesc}
			
		\item[spirit of revelation] \textbf{11} | \hyperref[ALMA-4-20]{Alma 4:20}; \hyperref[ALMA-5-46]{5:46}; \hyperref[ALMA-8-24]{8:24}; \hyperref[ALMA-9-21]{9:21}; \hyperref[ALMA-17-3]{17:3}; \hyperref[ALMA-23-6]{23:6}; \hyperref[ALMA-45-10]{45:10}; \hyperref[HELAMAN-4-23]{Helaman 4:23}; \hyperref[3NEPHI-3-19]{3 Nephi 3:19}; \hyperref[DC8-3]{DC 8:3}; \hyperref[DC11-25]{DC 11:25}
			\begin{compactdesc}
				% \item[Bible anchor verses] \phantom{.}
					% \begin{compactdesc}
						% \item[Romans 5:5] \gr{}
					% \end{compactdesc}
				\item[Commentary] What is the meaning of the word ``spirit'' in this phrase? It seems most natural to associate ``spirit'' here with the Holy Ghost, as thoug ``spirit'' is referring to the Holy Ghost's special function as communicator of revelation; the Holy Ghost appears in conjunction with the phrase in two places, Alma 5:46 and DC 8:3 (the reference to the Holy Ghost is in verse 2). Drawing from the \hyperref[SPIRIT-1828]{1828 dictionary entry} for ``spirit,'' though, there are other interesting meanings of ``spirit'' that could be in play here as well. I am especially interested in meaning \#11, ``Eager desire; disposition of mind excited and directed to a particular object.'' To make this meaning in the phrase more explicit, you'd have to change the wording a bit, but the idea (to use DC 8:3 as an example) would be something like: ``This is the state of mind which characterizes one who eagerly desire revelation; this is the state of mind Moses had when he brought the children of Israel through the Red Sea on dry ground.''
				
				No, this suggestion isn't perfect, and it doesn't completely fit with all the contexts in which the phrase occurs, but I think there is something important here about understanding the mindset that someone has who is actively seeking out revelation, or seeking to understand the revelation that the Holy Ghost is communicating to them throughout their day. They are eagerly desiring to live in the flow of the Spirit.
			\end{compactdesc}
			
		\item[spirit of truth] \textbf{16} | \hyperref[JOHN-14-17]{John 14:17}; \hyperref[JOHN-15-26]{15:26}; \hyperref[JOHN-16-13]{16:13}; \hyperref[1JOHN-4-6]{1 John 4:6}; \hyperref[DC6-15]{DC 6:15}; \hyperref[DC50-17]{DC 50:17} (2); \hyperref[DC50-19]{DC 50:19}; \hyperref[DC50-21]{DC 50:21} (2); \hyperref[DC93-9]{DC 93:9}; \hyperref[DC93-11]{DC 93:11}; \hyperref[DC93-23]{DC 93:23}; \hyperref[DC93-26]{DC 93:26} (2); \hyperref[DC107-71]{DC 107:71}
			% "spirit* of truth*"
			% \begin{compactdesc}
				% \item[Bible anchor verses] \phantom{.}
					% \begin{compactdesc}
						% \item[Romans 5:5] \gr{}
					% \end{compactdesc}
				% \item[Commentary] ---
			% \end{compactdesc}
			
	\end{compactdesc}
	
% % ----------------------
% \subsection{Bible Dictionaries} % *DICTIONARIES
% % ----------------------
	% \label{SPIRIT-BD}

% % ----------------------
% \subsection{Restoration Usage} % *RESTORATION
% % ----------------------
	% \label{SPIRIT-RESTORATION}

	% % -----
	% \subsubsection{Joseph Smith} % *JS
	% % -----
		% \label{SPIRIT-JS}
	
		% \begin{compactdesc}
			% \item[{\hyperref[KEY]{Date/Source}}]
		% \end{compactdesc}
	
	% % -----
	% \subsubsection{Brigham Young} % *BY
	% % -----
		% \label{SPIRIT-BY}
	
		% \begin{compactdesc}
			% \item[{\hyperref[KEY]{Date/Source}}]
		% \end{compactdesc}
	
% ----------------------
\subsection{Commentary and Studies} % *COMMENTARY *STUDIES
% ----------------------
	\label{SPIRIT-COMMENTARY}

	% % -----
	% \subsubsection{Key Points}
	% % -----
		% \label{SPIRIT-KEYS}
	
		% \begin{compactitem}
			% \item
		% \end{compactitem}
		
	% -----
	\subsubsection{The ``spirit of Christ''}
	% -----
		\label{SPIRIT-christ}
		
		See this idea at ``\hyperref[LIGHT-christ]{light of Christ}.''
		
	% -----
	\subsubsection{``Spirit'' as matter or something otherwise physical}
	% -----
		\label{SPIRIT-matter}
		
		\begin{compactitem}
			\item See the JSP/BC version of \hyperref[DC19-18]{DC 19:18}.
		\end{compactitem}
		
	% -----
	\subsubsection{On the possible equivalence of God's spirit, love, glory, etc.}
	% -----
		\label{SPIRIT-equivalence}
		
		See \hyperref[2NEPHI-4-21]{2 Nephi 4:21} and notes there.
		
		See also:
		
			\begin{compactitem}
				\item \hyperref[1NEPHI-11]{1 Nepih 11}, and the comments about God's love and how it ``sheds itself abroad.'' Nephi says it's the most desirable of all things, and the angel says it's the ``most joyous to the soul,'' so it somehow interacts with us to make us feel good things.
				\item \hyperref[1NEPHI-17-47]{1 Nephi 17:47--48}, full of the ``Spirit of God'' and the ``power of God.'' Verse 52 also.
				\item \hyperref[MOSIAH-13-5]{Mosiah 13:5--6}
			\end{compactitem} 